% Terjemahan Bahasa Indonesia dari prelude.tex baris 127-153.
% Sumber: Methods of Algebra, Volume 2, commit
% 9a5803ff77dd3257484cb177f851a73770a59dd3 (CC BY 4.0).
% stable-unit-id: o014.aljabr2.prelude.reader-guide

% segment-id: o014.li2.prelude.reader-guide.h001
\section*{Panduan Penggunaan}
\addcontentsline{toc}{section}{Panduan Penggunaan}
\hypertarget{o014.aljabr2.prelude.reader-guide}{ }

% segment-id: o014.li2.prelude.reader-guide.h002
\subsection*{Fungsi Buku}

% segment-id: o014.li2.prelude.reader-guide.p001
Untuk memenuhi fungsinya sebagai buku rujukan, buku ini perlu disusun cukup
sistematis. Karena itu, setiap bab terutama mengikuti
urutan logis. Hubungan ketergantungan antarbagiannya sangat kentara pada
\emph{Bagian Dalam}, sedangkan sejumlah kecil topik yang diperlukan tetapi
menyimpang dari jalur utama, atau yang lebih bersifat teknis, ditempatkan
dalam lampiran.

% segment-id: o014.li2.prelude.reader-guide.p002
Bagi pembelajar mandiri, bagian yang perlu dibaca bergantung pada minat dan
kebutuhan masing-masing. Pilihan utama terletak pada \emph{Bagian Luar}, sedangkan
\emph{Bagian Dalam} merupakan kerangka dasar bersama, meskipun pada skala subbab
masih ada sejumlah materi pilihan di dalamnya. Beberapa perincian yang
lebih abstrak atau teknis dapat dilewati seperlunya ketika pertama kali
belajar. Pembahasan yang bergantung pada lampiran sebaiknya dilewati dahulu atau
hanya dibaca sekilas. Saran terperinci mengenai pilihan-pilihan tersebut
terdapat pada penjelasan di awal setiap bab dan lampiran, terutama dalam
bagian \emph{Petunjuk Bacaan}.

% segment-id: o014.li2.prelude.reader-guide.p003
Seperti disebutkan sebelumnya ketika membahas tujuan buku, kita terpaksa
menyelubungi gagasan matematika dengan lapisan-lapisan teks.
Pembelajar mandiri harus mengerahkan upaya sungguh-sungguh untuk mengenali
matematika yang sesungguhnya di baliknya, layaknya menangkap kilat di
tengah gulungan awan tebal.

% segment-id: o014.li2.prelude.reader-guide.p004
Bagi pembaca yang menggunakan buku ini sebagai rujukan, penulis telah
berusaha menerapkan cara penulisan modular agar bagian-bagiannya mudah
dibaca secara terpisah. Hal ini tercermin bukan hanya dalam pembagian bab,
melainkan juga dalam dua hal berikut.
\begin{itemize}
	% segment-id: o014.li2.prelude.reader-guide.i001
	\item Sedapat mungkin, buku ini menggunakan notasi baku dalam pustaka
	kontemporer, disertai tinjauan ulang dan penjelasan. Bagian Konvensi pada
	akhir pendahuluan dan indeks di akhir buku merupakan dua sarana lain yang
	ampuh untuk memahami notasi.
	% segment-id: o014.li2.prelude.reader-guide.i002
	\item Berbagai pernyataan dan pembuktian dilengkapi dengan rujuk silang
	secara cermat, sehingga pembaca yang sudah memiliki pengetahuan tertentu
	tentang materinya dapat memahami uraian dengan cepat dan membangun peta
	konseptual secara efisien.
\end{itemize}

% segment-id: o014.li2.prelude.reader-guide.p005
Secara khusus, penulis berharap buku ini dapat membantu pembaca yang sedang
menyiapkan pengajaran atau mengikuti seminar untuk segera menghimpun
pengetahuan yang mereka perlukan, lalu mengolahnya dengan bahasa mereka sendiri.

% segment-id: o014.li2.prelude.reader-guide.p006
Mengenai penggunaan dalam pengajaran, harus diakui terus terang bahwa
penjadwalan perkuliahan tidak dapat diatur sesuka hati. Selama menulis
buku ini, penulis hampir tidak memperoleh kesempatan untuk menguji materi
ini di ruang kelas, sehingga hasilnya belum dapat dinilai. Berdasarkan
pengalaman terdahulu, sekalipun perincian dilewati di kelas, seluruh buku
ini jelas memuat bahan untuk lebih dari satu tahun akademik; \emph{Bagian Dalam}
mungkin dapat dipadatkan menjadi satu semester dengan penyajian ringkas. Jika
buku ini digunakan sebagai buku ajar utama, pengajar harus terlebih dahulu
membagi materi buku ini. Penulis berharap rancangan modular yang telah disebutkan
dapat mempermudah pekerjaan tersebut.

% segment-id: o014.li2.prelude.reader-guide.p007
Kurangnya pengujian di ruang kelas membawa dua akibat lain. Pertama, buku
ini tak pelak masih mengandung banyak kesalahan yang belum ditemukan.
Kedua, buku ini belum memiliki keluwesan dan kematangan khas buku ajar
klasik. Hal ini sebagian disebabkan oleh keadaan praktis dan sebagian lagi
oleh keterbatasan penulis; penulis hanya dapat memohon pengertian
pembaca.

% segment-id: o014.li2.prelude.reader-guide.h003
\subsection*{Tentang Latihan}

% segment-id: o014.li2.prelude.reader-guide.p008
Setiap bab dan lampiran diakhiri dengan latihan yang disusun mengikuti
urutan materi terkait dalam teks utama dan kadang-kadang disertai petunjuk.
Tujuan utamanya ialah menambahkan sifat-sifat atau hasil pelengkap,
menyediakan contoh, dan
memperluas wawasan. Sebagian latihan dimaksudkan untuk membantu pembaca
membiasakan diri dengan berbagai teknik dan pada umumnya tidak terlalu
teknis. Sejumlah kecil latihan meminta pembaca melengkapi perincian yang
dihilangkan dari teks utama; perincian itu kadang remeh, kadang mudah, dan
lebih sering keduanya sekaligus. Singkatnya, penulis berharap pembaca
pemula berusaha mengerjakan sebanyak mungkin latihan, tetapi tidak perlu
memaksakan diri atau merasa harus menyelesaikan semuanya.

% segment-id: o014.li2.prelude.reader-guide.h004
\subsection*{Kategori dan Ukurannya}

% segment-id: o014.li2.prelude.reader-guide.p009
Sebagian besar isi buku ini dinyatakan dalam bahasa kategori karena
praktik telah membuktikan bahwa teori kategori memang merupakan sarana
yang efisien untuk memahami dan menangani persoalan aljabar. Namun,
aljabar linear atau apa yang disebut aljabar homologis bukanlah cabang
teori kategori, sama seperti teori probabilitas bukanlah cabang teori
ukuran.

% segment-id: o014.li2.prelude.reader-guide.p010
Berbeda dari sejumlah buku ajar lain, buku ini dan Jilid I sama-sama
mensyaratkan agar kumpulan semua objek dan kumpulan semua morfisme suatu
kategori masing-masing merupakan himpunan; objek tidak boleh hanya
membentuk sebuah ``kelas''. Karena kelas sejati
(\emph{proper class}) tidak digunakan dalam definisi, secara formal kita harus memperkenalkan
konsep semesta Grothendieck (\emph{Grothendieck universe}) untuk mengukur
ukuran himpunan dan mengasumsikan bahwa semua himpunan termasuk dalam
suatu semesta $\mathcal{U}$; lihat
\cite[Asumsi~1.5.2, Catatan~1.5.6]{Li1}. Jika setiap persoalan dianalisis
secara tersendiri, aksioma mengenai keberadaan cukup banyak semesta---atau,
secara ekuivalen, cukup banyak kardinal tak terjangkau kuat
(\emph{strongly inaccessible cardinal})---sering kali dapat diperlemah.
Karena hal ini tidak menyangkut inti buku, kita tidak akan membahasnya
lebih jauh.
